\documentclass[]{article}
\usepackage{amsmath, amssymb, amsthm}

\usepackage{multicol, enumitem, physics, changepage}

\renewcommand{\thesubsection}{\thesection.\alph{subsection}}

\newcommand{\mrow}[3]{\left[
	\begin{array}{ccc}
		#1 & #2 & #3
	\end{array}
	\right]}
\newcommand{\mcol}[3]{\left[
	\begin{array}{c}
		#1\\
		#2\\
		#3
	\end{array}
	\right]}

\newenvironment{augmx}[1]{%
	\left[\begin{array}{@{}*{#1}{c}|c@{}}
	}{%
	\end{array}\right]
}

\newcommand{\extaugmx}[2]{
	\left[\begin{matrix}
		#1
	\end{matrix}\right|\left.\begin{matrix}
		#2
	\end{matrix}\right]
}

\newcommand{\xlrightarrow}{\xrightarrow{\hspace{0.8cm}}}
\newcommand{\reals}{\mathbb{R}}

\newcommand{\bmat}[1]{\begin{bmatrix}#1\end{bmatrix}}

\newcommand{\letmatrix}[2]{\mathcal{M}_{#1 \times #2}(\reals)}

\newcommand{\detmat}[1]{
	\left|\hspace{0.5em}
	\begin{matrix}
		#1
	\end{matrix}
	\hspace{0.5em}\right|
}

\begin{document}
\begin{center}
	{\Large\textbf{MATH545: Problem Set 5}}\\
	\large Grey Cole\\
	March 25, 2026
\end{center}
\section{}
Let $S_1$ be the set of all linear combinations of the vectors $\alpha_1, \alpha_2, ..., \alpha_k$ in $\reals^n$.\\\noindent
Let $S_2$ be the set of all linear combinations of the vectors $\alpha_1, \alpha_2, ..., \alpha_k, \alpha_1 + \alpha_2$.\\\noindent
Show that $S_1 = S_2$ by showing that each of the sets is a subset of the other.
\begin{proof}$\;$\\
	(1)  \begin{adjustwidth}{0.5cm}{}
	The definition of $S_2$ indicates $S_2 = S_1 \cup \alpha_1 + \vb{v_2}, \implies S_1 \subseteq S_2$.\\
	\end{adjustwidth}
	(2) \begin{adjustwidth}{0.5cm}{}
	By the definition of linear combinations,
	\[ S_1 = c_1\alpha_1 + c_2\alpha_2 + \cdots + c_k\alpha_k = \sum_{i=1}^{k}c_i\alpha_i \]
	Using $c_1 = 1, c_2 = 1,$ and $c_i = 0$ for all $i > 2$, 
	\[ 1\alpha_1 + 1\alpha_2 + 0\alpha_3 + \cdots + 0\alpha_k = \alpha_1 + \alpha_2 + 0[\alpha_3 + \cdots + \alpha_k] = \alpha_1 + \alpha_2 \]
	is necessarily in $S_1$.\\
	$\therefore \alpha_1 + \alpha_2 \in S_1$ which implies $S_1 \cup \alpha_1 + \alpha_2 \subseteq S_1$, therefore $S_2 \subseteq S_1$ by definition of $S_2$.
	\end{adjustwidth}
	Since we have shown $S_1 \subseteq S_2$ and $S_2 \subseteq S_1$, $S_1 = S_2$.
\end{proof}
\section{}
Suppose that $S = \{ \alpha_1, \alpha_2, \alpha_3 \}$ is linearly independent and further suppose that
\[ \beta_1 \alpha_2\;, \beta_2 = \alpha_1 + \alpha_3\;\text{, and } \beta_3 = \alpha_1 + \alpha_2 + \alpha_3\]
Determine whether the set $T = \{\beta_1, \beta_2, \beta_3\}$ is linearly dependent or linearly independent.\\\\

\textbf{Claim:} $T$ is linearly dependent.
\begin{proof}
	We will prove the claim by showing that there is a nontrivial solution to the equation
	\[ c_1\beta_1 + c_2\beta_2 + c_3\beta_3 = 0 \]
	i.e. a solution where at least one $c_1, c_2, c_3 \neq 0$. First, substitute all $\beta_i$ with the corresponding sum of vectors $\alpha$:
	\[ c_1\alpha_2 + c_2[\alpha_1 + \alpha_3] + c_3[\alpha_1 + \alpha_2 + \alpha_3] = 0 \]
	Let $c_3 \neq 0$ and $c_1 = c_2 = -c_3$, $\; \implies c_1,c_2,c_3 \neq 0$.
	\[ -c_3\alpha_2 - c_3[\alpha_1 + \alpha_3] + c_3[\alpha_1 + \alpha_2 + \alpha_3] = 0\]
	Reducing this equation, we get
	\[ -c_3[\alpha_1 + \alpha_2 + \alpha_3] + c_3[\alpha_1 + \alpha_2 + \alpha_3] = 0 \longrightarrow 0 = 0\]
	Thus we have found a nontrivial solution to the equation.\\
	$\therefore$ $T$ is linearly dependent.
\end{proof}
\section{}
\subsection{}
\begin{proof}
The set $\{M_1, M_2, M_3\}$ is linearly independent if there are no trivial solutions to the equation
\begin{equation}
	c_1\bmat{1&0\\-1&1} + c_2\bmat{1&1\\2&1} + c_3\bmat{0&1\\1&0} = \bmat{0&0\\0&0}
\end{equation}
	We can multiply each matrix by its scalar coefficient and solve the resulting linear system.
	\[ \bmat{c_1&0\\-c_1&c_1} + \bmat{c_2&c_2\\2c_2&c_2} + \bmat{0&c_3\\c_3&0} = \bmat{0&0\\0&0} = \bmat{c_1+c_2 & c_2+c_3 \\ -c_1 + 2c_2 + c_3&c_1 + c_2} \] 
	\begin{align}
		c_1 + c_2 = 0 = c_2 + c_3 \implies c_1 = c_3 = 0\\
		-(0) + 2c_2 + (0) = 0 = 2c_2 \implies c_2 = 0
	\end{align}
	Thus we have shown that $c_1 = c_2 = c_3 = 0$ which is necessary and sufficient to determine that the set $\{M_1, M_2, M_3\}$ is linearly independent.
\end{proof}
\subsection{}
\begin{equation}
	\label{3b.1}
	\bmat{1&-1\\2&1} = \bmat{c_1&0\\-c_1&c_1} + \bmat{c_2&c_2\\2c_2&c_2} + \bmat{0&c_3\\c_3&0} = \bmat{c_1 + c_2 & c_2 + c_3 \\ -c_1 + 2c_2 + c_3 & c_1 + c_2}
\end{equation}
From equation \ref{3b.1} we can discern the linear system
\begin{align*}
	c_1 + c_2 &= 1\\
	c_2 + c_3 &= -1\\
	-c_1 + 2c_2 + c_3 &= 2\\
	c_1 + c_2 &= 1
\end{align*}
We can exclude the duplicate equation and set up a matrix which we can row reduce and then use back substitution to find solutions to the system.
\begin{align*}
	\begin{augmx}{3}1&1&0&1\\0&1&1&-1\\-1&2&1&2\end{augmx} \longrightarrow \begin{augmx}{3}1&1&0&1\\0&1&1&-1\\0&0&-2&6\end{augmx} \longrightarrow \bmat{1&1&0\\0&1&1\\0&0&-2}\bmat{c_1\\c_2\\c_3} = \bmat{1\\-1\\6}
\end{align*}
\begin{align*}
	c_1 + c_2 = 1 & & c_3 = -3\\
	c_2 + c_3 = -1 & & c_2 = -1 + 3 \implies c_2 = 2\\
	-2c_3 = 6 & & c_1 = 1 - 2 \implies c_1 = -1
\end{align*}
Thus \boxed{c_1 = -1, c_2 = 2, c_3 = -3}.
\subsection{}
No, the matrix $\left[\begin{smallmatrix}1&-1\\1&2\end{smallmatrix}\right]$ cannot be written as a linear combination of $M_1, M_2,$ and $M_3$.
\begin{proof}
	Creating a linear system based on the possible linear combinations of $M_1, M_2, M_3$ creates a contradiction.
	\begin{equation}
		\bmat{1&-1\\1&2} = \bmat{c_1&0\\-c_1 & c_1} + \bmat{c_2&c_2\\2c_2&c_2} + \bmat{0&c_3\\c_3&0} = \bmat{c_1 + c_2&c_2 + c_3\\-c_1 + 2c_2 + c_3&c_1 + c_2}
	\end{equation}
	\begin{align*}
		c_1 + c_2 = 1\\
		c_2 + c_3 = -1\\
		-c_1 + 2c_2 + c_3 = 1\\
		c_1 + c_2 = 2
	\end{align*}
	$1 = c_1 + c_2 = 2$ implies $1 = 2$. This is a contradiction, and so there is no set of scalars $c_1,c_2,c_3$ that can create the given matrix from $M_1,M_2,M_3$.
\end{proof}
\subsection{}
\begin{equation}
	\label{3d.1}
	\bmat{a&b\\c&d} = \bmat{c_1 & 0\\-c_1 & c_1} + \bmat{c_2 & c_2\\2c_2 & c_2} + \bmat{0 & c_3\\c_3 & 0} = \bmat{c_1 + c_2 & c_2 + c_3\\-c_1 + 2c_2 + c_3 & c_1 + c_2}
\end{equation}
All matrices that can be written as a linear combination of $M_1, M_2$, and $M_3$ \emph{must} be able to be represented in terms of $c_1, c_2, c_3$ as described in equation \ref{3d.1}.\\
To put it another way, all matrices $\left[\begin{smallmatrix}a&b\\c&d\end{smallmatrix}\right]$ that can be written as a linear combination of $M_1, M_2$, and $M_3$ adhere to the following:
\begin{gather*}
	\left.
	\begin{aligned}
		c_1 + c_2 &= a = d \; \\
		c_2 + c_3 &= b\\
		-c_1 + 2c_2 + c_3 &= c
	\end{aligned}
	\right\}
	\begin{aligned}
		\\
		\;\; \text{such that }(c_1,c_2,c_3) \in \reals^3\\
		\;
	\end{aligned}
\end{gather*}
\begin{align*}
	\;
\end{align*}
\section{}
We can find a basis of $span(S)$ by row-reducing the matrix where each \emph{row} is an element in $S$.
\[\bmat{-2 & 0 & 2\\1 & 0 & -3\\-3&-3&-2\\1&2&-2} \longrightarrow \bmat{0 & 0 & 4\\1 & 0 & -3\\ 0 & -3 & 7\\ 0 & 2 & 1} \longrightarrow \bmat{0&0&1\\1&0&0\\0&-3&0\\0&2&0} \longrightarrow \bmat{0&0&1\\1&0&0\\0&1&0\\0&0&0}\]
It is trivial to see that the fourth row, corresponding to $\bmat{1\\2\\-2}$ in $S$, is not necessary to describe a vector in $span(S) \subseteq V = \reals^3$ (but the others are, as a basis of $span(S) \subseteq \reals^3$ must contain at least 3 vectors).\\ Thus, by removing the fourth element shown of $S$, a basis of $span(S)$ is \[ \boxed{\left\{\bmat{-2\\0\\2}, \bmat{1\\0\\-3}, \bmat{-3\\-3\\-2}\right\}} \]
\section{}
We can find a basis for $V = \reals^3$ which contains the given set $S$ as a subset by appending the standard basis of $\reals^3$ to $S$ and then removing those additions which are redundant (can be row-reduced to \textbf{0}).
\[S' = \left\{\bmat{2\\2\\-1}, \bmat{-1\\-1\\3}, \bmat{1\\0\\0}, \bmat{0\\1\\0}, \bmat{0\\0\\1}\right\}\]
\[ \bmat{2&2&-1\\-1&-1&3\\1&0&0\\0&1&0\\0&0&1} \rightarrow \bmat{0&0&1\\-1&-1&3\\1&0&0\\0&1&0\\0&0&1} \rightarrow \bmat{0&0&1\\0&-1&0\\1&0&0\\0&1&0\\0&0&0} \rightarrow \bmat{0&0&1\\0&1&0\\1&0&0\\0&0&0\\0&0&0}\]
Because the fourth and fifth row were eliminated, the fourth and fifth elements of $S'$ are not necessary to describe a vector in $V = \reals^3$. Therefore, a basis for $V = \reals^3$ which contains the given set as a subset is
\[ \left\{ \bmat{2\\2\\-1}, \bmat{-1\\-1\\3}, \bmat{1\\0\\0} \right\} \]
\end{document}